\chapter{SLG/Al$_{2}$O$_{3}$基板におけるコバルトインターカレーション}

本章ではグラフェンにおけるスピン軌道相互作用の増大を目指し，グラフェン/サファイア基板におけるコバルトのインターカレーションの研究についてまとめる．
まず化学蒸着(CVD:chemical vapor deposition)法によってSLGを作成した結果及びコバルトの蒸着について述べる．その後作成したSLGの質を評価した結果をまとめる．最後にコバルトインターカレーションについての結果を述べる．

\section{CVD法によるSLGの作成}
ここでは本研究で用いたSLGの作成について述べる．
グラフェンの作成方法にはいくつか候補がある．代表的なものはFig.\ref{fig:graphene_howtomake}にまとめた．

\begin{figure}[t]
 \begin{center}
  \includegraphics[width=120mm]{images/graphene_howtomake.eps}
  \end{center}
   \caption{グラフェンの代表的な作成方法．}
 \label{fig:graphene_howtomake}
\end{figure}


一つは初めてグラフェンの作成に成功したスコッチテープである．この方法はグラファイトの試料をスコッチテープを用いて剥離を繰り返すことで炭素の層を薄くしていきグラフェンを作成し基板に圧着する．この方法で作成したグラフェンは高品質であることが知られ，最大の移動度を測定されたものもスコッチテープ法で作られた．しかしスコッチテープで剥離を繰り返すため，手に入る試料のサイズは小さく加工も難しい．
SiCの熱分解を用いた方法はその名の通りSiC基盤を1000$\rm\ C^{\circ}$で加熱することにより表面にグラフェンが成長させる方法である．SiCを高温で加熱すると表面においてSiだけが選択的に脱離する．すると表面に残った炭素原子同士が自発的に結合しグラフェンを成長する．この方法によるグラフェンも質が良いことで知られる．しかし大面積のグラフェンを得られてもグラフェンの総数が制御しにくいことやコスト面が高価なためこの方法を選択しなかった．
化学気相蒸着，もしくは化学蒸着(CVD:Chemical Vipor Deposition)法は原料のガス（ここではエタノールガス）を流した中で基板を加熱することで基板表面で化学反応が起きグラフェンが成長する方法である．この方法は他の作成方法に比べグラフェンの質は劣るがコストが低く，基板の選択肢の多さや加工のしやすい．本研究ではこのCVD法を用いてグラフェンを作成した．




\section{ラマン分光法によるSLGの評価}

\section{コバルトインターカレーションの評価}

\subsection{XPS測定}

\subsection{AFM測定による評価}

